\chapter{The Idea of the Indian Village}

\section{The Village as an Ideological Category}

\begin{KeyConcept}
\begin{itemize}
\item The \textbf{idea of India} was that of an \textbf{agrarian, rural country} --- but this was a \textbf{Western (colonial) idea}. To a common-sense man, a village is a geographical space; to a sociologist, \textbf{the village is an ideological category through which the British conceptualised India} --- the \textbf{microcosm of the macrocosm}.
\item The British understood India through \textbf{ancient Sanskritic texts}, which described an agrarian, pastoral, religious, 'backward, uncivilized' India --- hence 'India is village and village is India'.
\item We know India (even the Bhagavad Gita) only through the British translations --- \textbf{colonial knowledge produced the knowledge of India}.
\end{itemize}
\end{KeyConcept}

\section{Colonial Discourse: Metcalfe, Marx and the Indologists}

\begin{KeyConcept}
\begin{itemize}
\item \textbf{Charles Metcalfe's} idea of India (the colonial/orientalist discourse --- to justify rule by first 'knowing' the ruled):
\item 1. \textbf{India consists of villages}.
\item 2. Villages are \textbf{little republics} (autonomous, politically sovereign units, with their own head).
\item 3. Villages are \textbf{self-sufficient} (economically autonomous --- everything produced, consumed and distributed within; no import/export).
\item 4. Villages are \textbf{static} and \textbf{resilient to change} (equilibrium maintained; 'dynasty after dynasty tumbled, revolution succeeded revolution, but village India does not reflect any change').
\item 5. \textbf{Communal ownership of land}.
\item 6. \textbf{Functional integration of various occupational groups}.
\end{itemize}
\end{KeyConcept}

\begin{KeyConcept}
\begin{itemize}
\item \textbf{Marx} (the \textbf{Asiatic mode of production}) described India the same way --- self-sufficient little republics, no private ownership, communal land, integrated occupational groups. (A. R. Desai, as a Marxist, reproduced this view.)
\item An \textbf{Indologist} described the Indian village as: a \textbf{traditional community}, \textbf{self-sufficient}, \textbf{patriarchal in governance}, and \textbf{surrounded by hostile villages} (hence autonomous, under a despotic government outside).
\item The first historian of Indian civilisation, \textbf{James Mill}, divided Indian history into \textbf{ancient (Hindu), medieval (Muslim) and modern (British)} --- constructing the common-sense we still inherit (even \textbf{Bipan Chandra} reproduces this colonial version of history, which is why the subaltern and tribal movements go missing).
\end{itemize}
\end{KeyConcept}

\section{Nationalist Discourse: Gandhi, Nehru, Ambedkar}

\begin{KeyConcept}
\begin{itemize}
\item \textbf{Gandhi} --- \textbf{'the real India lies in the villages (Bharat)'}: the village represents the \textbf{essential core of Indian civilisation} and is the \textbf{site of authenticity} ('India begins and ends in the village'). Develop India by \textbf{developing the village} --- he was against \textbf{westernized development} (which replaces tradition), not development per se. \textbf{Real Swaraj (self-rule) = revival of village communities} --- hence \textbf{Panchayati Raj} (a Gandhian principle, in the DPSPs). Gandhi was \textbf{against the idea of the state}.
\item \textbf{Nehru} --- the \textbf{village is traditional and backward India}, a \textbf{den of superstition, tradition and underdevelopment} --- which cannot gel with a modern industrial society. Instead of reviving the village, he suggested \textbf{modern industrial development} (heavy machinery, technology, dams) and modernising \textbf{agriculture through modern technology}. India should be a \textbf{strong nation state} (villages should not be autonomous like princely states).
\item \textbf{Ambedkar} --- the \textbf{Dalitist idea of the village}: the village is a \textbf{den of oppression, violence, exclusion and corruption} --- a \textbf{Hindu Republic which is the negation of the Republic}, a colonialism of the Hindus designed to oppress the Dalits. The village is the \textbf{'working plant of the Hindu caste order'}; the real India is not yet free.
\end{itemize}
\end{KeyConcept}

\begin{CriticalPoint}
Even field studies reproduced the upper-caste view: \textbf{Beteille}, to befriend the villagers of Shripuram, used his \textbf{Brahmin identity} (not his professor/researcher identity); the Brahmin headman (a landlord) helped him, on the condition that he would \textbf{never visit Cheri} (the untouchable cluster). So whatever Beteille recorded was the \textbf{upper-caste view of Shripuram} --- just as we read the upper-caste view of modern Indian history.
\end{CriticalPoint}

\section{The State View: The Village as Underdeveloped}

\begin{KeyConcept}
\begin{itemize}
\item The \textbf{Indian state (Nehru's vision)} looked at the village as \textbf{underdeveloped} --- through the \textbf{Western idea of modernization} (Rostow's five stages: traditional $\rightarrow$ pre-take-off $\rightarrow$ take-off $\rightarrow$ drive to maturity $\rightarrow$ age of mass consumption).
\item It established the \textbf{Planning Commission} (economists and social anthropologists, many from the London School of Economics, who hired unpaid student interns to gather village data).
\item Programs: \textbf{IRDP (Integrated Rural Development Programme), CDP (Community Development Programme), land reforms, the Green Revolution}.
\item \textbf{Economists} studied only the economic aspects (production, distribution, supply chain); \textbf{anthropologists} studied only the socio-cultural aspects (everyday life, rituals, dress, food, festivals). \textbf{M. N. Srinivas} (in \textbf{'India's Villages'}) observed that this \textbf{mutual exclusion of aspects distorted our comprehension of the rural social structure}.
\end{itemize}
\end{KeyConcept}

\begin{CriticalPoint}
\begin{itemize}
\item \textbf{One modernity was imposed on multiple traditions}: multiple villages, but one plan for all. This \textbf{obliterated the diversity and specificity} of villages and reduced them to a \textbf{monolithic village} dependent on state policy.
\item \textbf{Jan Breman}: the village was \textbf{nationalized, modernized and anthropologized}.
\item \textbf{Thakur}: the village became a \textbf{litmus test} for the presence, absence or degree of development; focusing only on underdevelopment encourages a \textbf{unified, monolithic village India} dependent on the state --- a reproduction of colonialism (sociologists became 'insider-outsiders' treating the insider as outsider).
\end{itemize}
\end{CriticalPoint}

\section{Village Studies: Structural-Functional and Marxist}

\begin{KeyConcept}
\begin{itemize}
\item \textbf{Village studies} in India were initiated by the \textbf{Planning Commission} in the \textbf{1950s--60s}.
\item \textbf{Structural-functional branch}: \textbf{William Wiser} (Karimpur), \textbf{M. N. Srinivas} (Rampura), \textbf{S. C. Dube} (Shamirpet).
\item \textbf{Marxist branch (from the 1970s)}: \textbf{A. R. Desai} (planning, cooperatives and land reforms actually favoured \textbf{bourgeois industrialization} --- socialism was only paid lip service, the instruments were 'socialist' but controlled by capital) and \textbf{Ramkrishna Mukherjee} (who established the base for a Marxist study of the village).
\end{itemize}
\end{KeyConcept}

\section{William Wiser and the Jajmani System}

\begin{KeyConcept}
\begin{itemize}
\item \textbf{William Wiser} (Karimpur) was the \textbf{first to study the jajmani relationship} in the field.
\item \textbf{Jajmani} is a relationship between an upper and a lower caste --- e.g. the \textbf{nai (barber)} cuts the pandit's hair for free, and the Brahmin in exchange gives grains --- \textbf{service for goods, no cash} (it is \textbf{not barter}).
\item Wiser's conclusion: \textbf{nobody is master and nobody is servant} --- the one who provides becomes the servant of the one provided to, but the latter also provides something in return. The relation is one of \textbf{interdependence (harmony)} --- a structural-functional view.
\end{itemize}
\end{KeyConcept}

\section{S. C. Dube and the Dominant Individual}

\begin{KeyConcept}
\begin{itemize}
\item \textbf{S. C. Dube} (Shamirpet, Telangana) critiqued Srinivas's 'dominant caste': it is \textbf{not the caste but the dominant individuals/families} that are dominant --- the whole village depends on them.
\item He identified \textbf{six features of a dominant individual}: (1) \textbf{individual traits/charisma}, (2) \textbf{seniority (age)}, (3) \textbf{wealth}, (4) \textbf{numbers}, (5) \textbf{education}, (6) \textbf{a government job}.
\item Such charismatic, senior, wealthy, educated figures (e.g. a leader like Sonia Gandhi in Rae Bareli, or Mulayam Singh Yadav) neutralise conflict --- \textbf{harmony is maintained} --- so there is little scope for conflict in the village.
\end{itemize}
\end{KeyConcept}

\begin{QuickRevision}
\begin{itemize}
\item Village = ideological category (microcosm of macrocosm) through which British conceptualised India.
\item Colonial discourse: Metcalfe (villages = little republics, self-sufficient, static, communal land, functional integration); Marx (Asiatic mode of production); Indologist (traditional, self-sufficient, patriarchal, surrounded by hostile villages); James Mill (ancient/medieval/modern).
\item Nationalist discourse: Gandhi (real India in villages; Swaraj = revival; Panchayati Raj), Nehru (village backward; modern industry), Ambedkar (village = den of oppression; Hindu Republic = negation of Republic; working plant of caste order).
\item Beteille's Brahmin identity / Cheri exclusion = upper-caste view reproduced.
\item State view: village underdeveloped (modernization theory); Planning Commission; IRDP/CDP/land reform/Green Revolution; economists vs anthropologists (Srinivas's 'India's Villages'); one modernity on multiple traditions; monolithic village; Jan Breman; Thakur (litmus test).
\item Village studies (1950s-60s): structural-functional (Wiser/Karimpur/jajmani, Srinivas/Rampura, Dube/Shamirpet) and Marxist (A. R. Desai, Ramkrishna Mukherjee).
\item Jajmani = service-for-goods, no cash, interdependence/harmony (not master-servant).
\item Dube's dominant individual: charisma, age, wealth, numbers, education, government job; conflict neutralized.
\end{itemize}
\end{QuickRevision}

